% TOP_PROBLEM: {"id": 4800011, "title": "Multiple ergodic averages — Problem 11", "category_id": 11, "category": {"id": 11, "name": "dynamical_systems", "display_name": "Dynamical Systems", "description": "Problems about long-term behavior of deterministic systems, Hamiltonian dynamics, and periodic orbits.", "slug": "dynamical-systems", "order_index": 11, "created_at": "2026-07-31T15:26:25.671Z"}, "status": "partially_solved", "problem_number": "AMR-047-0011", "set_id": 13, "statusQualification": "Retains the triple-linear part; the n,n^2 bilinear part is explicitly identified as solved."}
% FIRST_POSTED: 2026-10-01
% ENTRY_KIND: statement_only
% UnsolvedMath ID 4800011; source catalogue code AMR-047-0011
% !TeX program = lualatex
% Definition and sources only; this file is not an LLM solution attempt.
\documentclass[11pt,a4paper]{article}
\usepackage[margin=25mm]{geometry}
\usepackage{fontspec}
\setmainfont{lmroman10-regular.otf}[BoldFont=lmroman10-bold.otf,
 ItalicFont=lmroman10-italic.otf,BoldItalicFont=lmroman10-bolditalic.otf]
\usepackage{amsmath,amssymb,mathtools}
\usepackage{unicode-math}
\setmathfont{latinmodern-math.otf}
\AtBeginDocument{\let\setminus\smallsetminus}
\usepackage{xurl,hyperref}
\hypersetup{colorlinks=true,urlcolor=blue}
\setcounter{secnumdepth}{0}
\setlength{\parindent}{0pt}
\setlength{\parskip}{5pt}
\setlength{\emergencystretch}{3em}
\begin{document}
\section*{Multiple ergodic averages — Problem 11}\label{4800011}
{\small UnsolvedMath ID 4800011 · AMR-047-0011 · Dynamical Systems · AMR Open Problem Lists}
\subsection{Definitions and mathematical statement}
Let $(X,\mathcal B,\mu)$ be a probability space and let
$T:X\to X$ be invertible and measure preserving. For
essentially bounded measurable complex functions $f,g,h$,
form the nonconventional averages
\[
 A_N(x)=\frac1N\sum_{n=1}^N
          f(T^n x)g(T^{2n}x)h(T^{3n}x).
\]
The retained target is to prove that for every such
system and every such triple, the finite complex limit
$\lim_{N\to\infty}A_N(x)$ exists for $\mu$-almost every
$x$. The exceptional null set may depend on the system
and on $f,g,h$. No ergodicity or mixing assumption is
imposed, and no formula for the limit is prescribed.

This asks for convergence at almost every point along
the full sequence of averaging lengths. Convergence in
$L^2(\mu)$, convergence in measure, or convergence on a
subsequence would not establish that conclusion. Choices
of representatives of the functions are immaterial after
removing the countable union of their iterate-preimages
of null sets.

The original numbered problem also asks for pointwise
convergence of the bilinear averages at times $n$ and $n^2$.
That second assertion was proved by Krause, Mirek, and
Tao, as recorded in the author's progress update. It is
included here only to specify the scope of the source
entry; the new unresolved target is its triple linear
average assertion.
\subsection{Short English statement}
Prove almost-everywhere convergence of the triple averages at times n, 2n, and 3n.
\subsection{Sources}
\begin{itemize}
\item[S1] ULAM.AI and the UnsolvedMath Contributors, supplied problems.json, record 4800011 (AMR-047-0011), AMR Open Problem Lists. Catalogue identity and source transcription; CC BY 4.0.
\url{https://huggingface.co/datasets/ulamai/UnsolvedMath}.
\item[S2] Frantzikinakis - Some open problems on multiple ergodic averages (2016), Problem 11. Original source indexed by AMR-047-0011.
\url{https://arxiv.org/abs/1103.3808}.
\item[S3] Frantzikinakis, November 2025 progress update for Problem 11.
\url{https://users.math.uoc.gr/~frantzikinakis/OpenProblemsNew.html}.
\item[S4] Krause, Mirek and Tao, Pointwise ergodic theorems for non-conventional bilinear polynomial averages.
\url{https://doi.org/10.4007/annals.2022.195.3.3}.
\end{itemize}
\end{document}
